\documentclass[a4paper, 12pt]{article}

%%% BuS LAB PROTOCOLL PREAMBLE
%%% 2019
%%%%%%%%%%%%%%%%%%%%%%%%%%%%%%%


%%% PACKAGES
%%%%%%%%%%%%%%%%%%%%%%%%%%%

\usepackage[ngerman]{babel}

\usepackage[utf8]{inputenc}
\usepackage{amsmath}
\usepackage{pgfplots}
\usepackage{tikz}
\usepackage[many]{tcolorbox}
\usepackage{graphicx,import}
\graphicspath{ {./graphics/} }
\usepackage{pdfpages}
\usepackage{dashrule}
\usepackage{float}
\usepackage{siunitx}
\usepackage{booktabs}
\usepackage{transparent}
\usepackage{svg}
\usepackage[version=4]{mhchem}
\usepackage[bottom]{footmisc}
\usepackage[european, straightvoltages]{circuitikz}
\usepackage{multicol}


%%% DOCUMENT GEOMETRY
%%%%%%%%%%%%%%%%%%%%%%%%%%%

\usepackage{geometry}
\geometry{
 a4paper,
 total={0.6180339887498948\paperwidth,0.6180339887498948\paperheight},
 top = 0.1458980337503154\paperheight,
 bottom = 0.1458980337503154\paperheight
 }
\setlength{\jot}{0.013155617496424828\paperheight}
\linespread{1.1458980337503154}

\setlength{\parskip}{0.013155617496424828\paperheight} % paragraph spacing


%%% COLORS
%%%%%%%%%%%%%%%%%%%%%%%%%%%

\definecolor{red1}{HTML}{f38181}
\definecolor{yellow1}{HTML}{fce38a}
\definecolor{green1}{HTML}{95e1d3}
\definecolor{blue1}{HTML}{66bfbf}
\definecolor{hsblue}{HTML}{00b1db}
\definecolor{hsgrey}{HTML}{afafaf}

%%% CONSTANTS
%%%%%%%%%%%%%%%%%%%%%%%%%%%
\newlength{\smallvert}
\setlength{\smallvert}{0.0131556\paperheight}


%%% COMMANDS
%%%%%%%%%%%%%%%%%%%%%%%%%%%

% differential d
\newcommand*\dif{\mathop{}\!\mathrm{d}}

% horizontal line
\newcommand{\holine}[1]{
  	\begin{center}
	  	\noindent{\color{hsgrey}\hdashrule[0ex]{#1}{1pt}{3mm}}\\%[0.0131556\paperheight]
  	\end{center}
}

% mini section
\newcommand{\minisec}[1]{ \noindent\underline{\textit {#1} } \\}

% quick function plot
\newcommand{\plotfun}[3]{
  \vspace{0.021286\paperheight}
  \begin{center}
    \begin{tikzpicture}
      \begin{axis}[
        axis x line=center,
        axis y line=center,
        ]
        \addplot[draw=red1][domain=#2:#3]{#1};
      \end{axis}
    \end{tikzpicture}
  \end{center}
}

% box for notes
\newcommand{\notebox}[1]{

\tcbset{colback=white,colframe=red1!100!black,title=Note!,width=0.618\paperwidth,arc=0pt}

 \begin{center}
  \begin{tcolorbox}[]
   #1 
  \end{tcolorbox}
 
 \end{center} 
 
}

% box for equation
\newcommand{\eqbox}[2]{
	
	\tcbset{colback=white,colframe=hsblue!100!black,title=,width=#2,arc=0pt}
	
	\begin{center}
		\begin{tcolorbox}[ams align*]
				#1
		\end{tcolorbox}
		
	\end{center} 
	
}

\newcommand{\inklude}[1]{%
    \def\svgwidth{\columnwidth}
    \import{./graphics/}{#1.pdf_tex}
}
% END OF PREAMBLE

%%%%%%%%%%%%%%%%%%%%%%%%%%%%%%%%%%%%%

\begin{document}

%%%%%%%%%%%%%%%%%%%%%%%%%%%%%%%%%%%%%
  \includepdf{./titlepage/titlepage.pdf}
  \clearpage
  \setcounter{page}{1}
%%%%%%%%%%%%%%%%%%%%%%%%%%%%%%%%%%%%%

\section{Konfiguration des MSP430 zur Port-Ein- und Ausgabe}
Jeder Port besteht aus 8 physischen Pins (Px.0 bis Px.7) und hat entsprechend mehrere 8-bit Register, die zur
Konfiguration dienen. Die für die Ein- und Ausgabefunktion relevanten Register
sind\footnote{x ist die Portnummer}:
\begin{enumerate}
\item[] \textbf{PxSEL}
  \begin{enumerate}
    \item[] \textit{Function Select Register}: Legt fest, welche Funktion verwendet wird.
      \begin{enumerate}
        \item[0:]  I/O-Funktion
        \item[1:]  Peripherie-Funktion
      \end{enumerate}
  \end{enumerate}
\item[] \textbf{PxDIR}
  \begin{enumerate}
  \item[] \textit{Direction Register}: Legt die Richtung, also Ein- oder Ausgabe fest.
  \begin{enumerate}
    \item[0:]  Input
    \item[1:]  Output
  \end{enumerate}
  \end{enumerate}
\item[] \textbf{PxOUT/PxIN}
  \begin{enumerate}
\item[]    \textit{Input/Output Register}: Bestimmen, abhängig vom Wert des PxDIR, den
   eigentlichen Ein- bzw. Ausgabewert.
 \end{enumerate}
\item[] \textbf{PxREN}
  \begin{enumerate}
    \item[] \emph{Resistor ENable Register}: Legt fest, ob Pullup- oder Pulldown
      Widerstände verwendet werden sollen (PxDIR auf Input).
      \begin{enumerate}
        \item[0:] Widerstand deaktiviert
        \item[1:] Der Wert in \emph{PxOUT} entscheidet ob Pullup/down
          \begin{enumerate}
            \item[0:] Pulldown
            \item[1:] Pullup
          \end{enumerate}
      \end{enumerate}
  \end{enumerate}
\end{enumerate}

Beispiel: Ausgabe HIGH an P1.0
\vspace{\parskip}
\begin{lstlisting}
P1SEL &=~ BIT0; // lösche Bit0, select IO 
P1DIR |=  BIT0; // Direction: Output
P1OUT |=  BIT0; // 1 an P1.0
\end{lstlisting}


\section{I/O-Funktionen}

\section{Aufgaben der Module des Rahmenprogramms}

  
 \end{document}

% Local Variables:
% TeX-engine: luatex
% End:
